\documentclass[12pt,a4paper]{article}

% ─── Packages ────────────────────────────────────────────────────────────────
\usepackage[utf8]{inputenc}
\usepackage[T1]{fontenc}
\usepackage[french]{babel}
\usepackage{geometry}
\geometry{margin=2.5cm}

\usepackage{amsmath, amssymb, amsthm}
\usepackage{algorithm}
\usepackage{algpseudocode}
\usepackage{tikz}
\usetikzlibrary{arrows.meta, positioning, shapes.geometric, decorations.markings, fit, backgrounds}
\usepackage{pgfplots}
\pgfplotsset{compat=1.18}
\usepackage{booktabs}
\usepackage{array}
\usepackage{xcolor}
\usepackage{hyperref}
\usepackage{enumitem}
\usepackage{mdframed}
\usepackage{lmodern}
\usepackage{microtype}
\usepackage{graphicx}
\usepackage{caption}
\usepackage{tcolorbox}
\tcbuselibrary{skins, breakable}

% ─── Colors ──────────────────────────────────────────────────────────────────
\definecolor{urgence}{RGB}{180,30,30}
\definecolor{theorie}{RGB}{30,80,160}
\definecolor{pratique}{RGB}{20,130,80}
\definecolor{attention}{RGB}{200,120,0}
\definecolor{grisléger}{RGB}{245,245,245}
\definecolor{bleuléger}{RGB}{230,240,255}
\definecolor{vertléger}{RGB}{230,248,238}
\definecolor{twist}{RGB}{110,30,160}
\definecolor{twistléger}{RGB}{245,230,255}

% ─── Custom environments ──────────────────────────────────────────────────────
\newtcolorbox{defbox}[1]{
  colback=bleuléger, colframe=theorie, fonttitle=\bfseries,
  title={#1}, breakable, sharp corners=south
}

\newtcolorbox{metierbox}[1]{
  colback=vertléger, colframe=pratique, fonttitle=\bfseries,
  title={#1}, breakable, sharp corners=south
}

\newtcolorbox{alertbox}[1]{
  colback=yellow!10, colframe=attention, fonttitle=\bfseries,
  title={#1}, breakable, sharp corners=south
}

\newtcolorbox{jurybox}[1]{
  colback=red!5, colframe=urgence, fonttitle=\bfseries,
  title={#1}, breakable, sharp corners=south
}

\newtcolorbox{twistbox}[1]{
  colback=twistléger, colframe=twist, fonttitle=\bfseries\large,
  title={#1}, breakable, sharp corners=south
}

% ─── Theorem styles ───────────────────────────────────────────────────────────
\theoremstyle{definition}
\newtheorem{definition}{Définition}[section]
\newtheorem{propriete}{Propriété}[section]
\newtheorem{invariant}{Invariant métier}[section]

\theoremstyle{remark}
\newtheorem{remarque}{Remarque}[section]
\newtheorem{exemple}{Exemple}[section]

% ─── Hyperref setup ───────────────────────────────────────────────────────────
\hypersetup{
  colorlinks=true,
  linkcolor=theorie,
  urlcolor=urgence,
  pdftitle={Twist 5 — Saturation du destinataire},
  pdfauthor={Système de routage critique},
}

% ─────────────────────────────────────────────────────────────────────────────
\begin{document}

% ═══════════════════════════════════════════════════════════════════
%  COUVERTURE
% ═══════════════════════════════════════════════════════════════════
\begin{titlepage}
  \centering
  \vspace*{2cm}
  {\color{twist}\rule{\linewidth}{2pt}}\\[0.5em]
  {\Huge\bfseries Twist 5}\\[0.4em]
  {\Large\bfseries L'hôpital le plus proche\\[0.2em]
  n'accepte plus de nouveaux cas de trauma}\\[0.4em]
  {\color{twist}\rule{\linewidth}{2pt}}\\[2em]

  % Schéma illustratif : deux hôpitaux, le plus proche barré
  \begin{tikzpicture}
    \node[circle, fill=urgence!20, draw=urgence, thick, minimum size=1.8cm,
          font=\small\bfseries] (amb) at (0,0) {AMB};

    % Hôpital proche — SATURÉ
    \node[rectangle, rounded corners=4pt, fill=urgence!10, draw=urgence,
          thick, minimum width=2.4cm, minimum height=1.2cm,
          font=\small\bfseries] (h1) at (4.5, 1.8) {CHU};
    \node[font=\footnotesize\bfseries, color=urgence] at (4.5, 0.95)
          {SATURÉ};
    \draw[-{Latex[length=3mm]}, urgence, thick, dashed]
          (amb) -- node[above,sloped,font=\small]{$d=9\,\text{min}$} (h1);
    % croix rouge sur le CHU
    \draw[urgence, very thick, line cap=round]
          (3.35,1.15) -- (5.65,2.45);
    \draw[urgence, very thick, line cap=round]
          (3.35,2.45) -- (5.65,1.15);

    % Hôpital alternatif — DISPONIBLE
    \node[rectangle, rounded corners=4pt, fill=pratique!10, draw=pratique,
          thick, minimum width=2.8cm, minimum height=1.2cm,
          font=\small\bfseries] (h2) at (5.5,-1.8) {Hôpital Central};
    \node[font=\footnotesize\bfseries, color=pratique] at (5.5,-2.6)
          {disponible};
    \draw[-{Latex[length=3mm]}, pratique, very thick]
          (amb) -- node[below,sloped,font=\small]{$d=14\,\text{min}$} (h2);

    % Légende
    \node[font=\footnotesize, color=urgence] at (2.25,  3.0)
          {Destination initiale — refusée};
    \node[font=\footnotesize, color=pratique] at (3.5, -3.2)
          {Destination recalculée — acceptée};
  \end{tikzpicture}\\[2em]

  {\large Note technique — Gestion de la saturation hospitalière}\\[1em]
  {\large\color{gray} Extension du moteur de routage dynamique A*}\\[3em]

  \begin{tabular}{rl}
    \textbf{Sujet :}    & Sujet 07 — Hackverse 2026 \\
    \textbf{Twist :}    & \#5 — Saturation du destinataire \\
    \textbf{Impact :}   & Destination devient une variable, non un paramètre fixe \\
    \textbf{Véhicule :} & Ambulance (cas principal), applicable aux autres types \\
  \end{tabular}

  \vfill
  {\small\color{gray} Document pédagogique — Présentation jury et équipe technique}
\end{titlepage}

% ═══════════════════════════════════════════════════════════════════
\tableofcontents
\newpage

% ═══════════════════════════════════════════════════════════════════
\section*{Introduction}
\addcontentsline{toc}{section}{Introduction}

Jusqu'ici, le moteur de routage traitait la \textbf{destination comme une
constante} : l'opérateur saisit les coordonnées d'un hôpital, l'algorithme
A* calcule le meilleur chemin pour y arriver. La seule complexité était
celle du réseau routier — trafic, météo, passabilité.

\begin{twistbox}{Twist 5 — Le problème change de nature}
L'hôpital le plus proche \textbf{n'accepte plus} de nouveaux cas de trauma.
Ce n'est pas une alerte sur une route : c'est une alerte sur la
\emph{destination elle-même}. Le système ne peut plus supposer que la
destination est valide. Il doit désormais \textbf{choisir où aller}
avant même de calculer comment y aller.
\end{twistbox}

Ce twist introduit une nouvelle dimension de complexité : la destination
est maintenant une \textbf{variable dynamique} au même titre que le poids
d'un arc. Il faut modéliser la disponibilité des hôpitaux, sélectionner
le meilleur destinataire parmi les alternatives, et recalculer si la
situation change en cours de route.

\newpage

% ═══════════════════════════════════════════════════════════════════
\section{Impact sur le modèle existant}

\subsection{Ce qui changeait dans les twists précédents}

Dans les itérations précédentes du moteur, tous les changements affectaient
les \textbf{arcs du graphe} : congestion, pluie, route bloquée, passabilité
selon le type de véhicule. La destination restait un nœud fixe, connu
à l'avance, toujours valide.

\subsection{Ce que le twist 5 modifie}

Le twist 5 rend les \textbf{nœuds de destination} potentiellement invalides.
Cela a trois conséquences directes sur l'architecture :

\begin{enumerate}[leftmargin=2em, itemsep=6pt]
  \item \textbf{Les hôpitaux ne sont plus de simples nœuds du graphe.}
    Ils deviennent des entités avec un état dynamique : capacité courante,
    statut d'accueil, spécialités disponibles.

  \item \textbf{L'algorithme de routage doit être précédé d'une étape de
    sélection de destination.} On ne peut plus appeler A*$(s, t)$ avec
    $t$ fixe — il faut d'abord déterminer quel $t \in \mathcal{H}$ est
    valide et optimal.

  \item \textbf{Le recalcul dynamique peut être déclenché par un événement
    hospitalier} (saturation en cours de trajet), et pas seulement par
    un événement routier.
\end{enumerate}

\begin{alertbox}{Ce que le système ne peut plus supposer}
\begin{itemize}[leftmargin=1.5em]
  \item \textbf{Interdit :} « La destination est toujours l'hôpital le
    plus proche. »
  \item \textbf{Interdit :} « Si l'ambulance est en route, l'hôpital de
    destination sera encore disponible à l'arrivée. »
  \item \textbf{Requis :} interroger le statut des hôpitaux \emph{avant}
    de calculer l'itinéraire \emph{et} surveiller ce statut pendant le trajet.
\end{itemize}
\end{alertbox}

\newpage

% ═══════════════════════════════════════════════════════════════════
\section{Modélisation : les hôpitaux comme nœuds dynamiques}

\subsection{Extension du graphe routier}

Le graphe routier existant $G = (V, E, w)$ est étendu avec un sous-ensemble
de nœuds hospitaliers $\mathcal{H} \subset V$, auxquels on attache un
vecteur d'état dynamique.

\begin{defbox}{Définition : Nœud hospitalier avec état dynamique}
Un nœud hospitalier $h \in \mathcal{H}$ est un nœud $v \in V$ enrichi
d'un vecteur d'état :
\[
  \sigma(h, t) =
  \bigl(\,
    \texttt{accepts\_trauma}(h,t),\;\;
    \texttt{capacity}(h,t),\;\;
    \texttt{eta\_reopening}(h,t)
  \,\bigr)
\]
où :
\begin{itemize}[leftmargin=1.5em]
  \item $\texttt{accepts\_trauma}(h,t) \in \{0, 1\}$ — vaut $1$ si l'hôpital
    accepte de nouveaux cas de trauma à l'instant $t$ ;
  \item $\texttt{capacity}(h,t) \in [0, 1]$ — taux de remplissage (0 = vide,
    1 = saturé) ;
  \item $\texttt{eta\_reopening}(h,t) \in \mathbb{R}^+ \cup \{+\infty\}$ —
    durée estimée avant réouverture (en minutes), ou $+\infty$ si inconnue.
\end{itemize}
\end{defbox}

\subsection{Ensemble des destinations candidates}

À chaque instant $t$, l'ensemble des destinations valides pour une
ambulance transportant un cas de trauma est :

\[
  \mathcal{H}^*(t) = \bigl\{\, h \in \mathcal{H}
    \;\mid\; \texttt{accepts\_trauma}(h,t) = 1 \,\bigr\}
\]

Si $\mathcal{H}^*(t) = \emptyset$ — c'est-à-dire qu'aucun hôpital
de la zone ne peut accepter le patient — le système doit déclencher
une procédure d'escalade (voir section~\ref{sec:limites}).

\subsection{Fonction de coût multi-destination}

Le problème de routage avec destination variable peut se formuler comme :

\[
  h^* = \arg\min_{h \in \mathcal{H}^*(t)}
  \Bigl[\,
    d_G(s, h)
    \;+\;
    \alpha \cdot \texttt{capacity}(h, t)
    \;+\;
    \beta \cdot \mathbb{1}[\texttt{capacity}(h,t) > \tau]
  \,\Bigr]
\]

où :
\begin{itemize}[leftmargin=1.5em]
  \item $d_G(s, h)$ est le temps de trajet optimal de la source $s$
    à l'hôpital $h$ dans le graphe dynamique ;
  \item $\alpha \cdot \texttt{capacity}(h, t)$ est une pénalité
    proportionnelle au taux de remplissage — préférer les hôpitaux
    moins encombrés, toutes choses égales par ailleurs ;
  \item $\beta \cdot \mathbb{1}[\texttt{capacity}(h,t) > \tau]$ est
    une pénalité binaire si le taux dépasse un seuil $\tau$ (typiquement
    $\tau = 0{,}85$), signalant un risque imminent de saturation avant
    l'arrivée du véhicule.
\end{itemize}

Les paramètres $\alpha$ et $\beta$ sont configurables par l'opérateur
selon la politique de soin de la zone. En mode par défaut :
$\alpha = 2\,\text{min}$ et $\beta = 5\,\text{min}$.

\begin{metierbox}{Interprétation opérationnelle}
Le coût total n'est pas simplement « le trajet le plus court » mais
« le trajet le plus court vers un hôpital qui pourra réellement
accueillir le patient à l'arrivée ». Un hôpital à 14 minutes à
moitié vide peut être préférable à un hôpital à 12 minutes presque
saturé, si le risque qu'il ferme pendant le trajet est non négligeable.
\end{metierbox}

\newpage

% ═══════════════════════════════════════════════════════════════════
\section{Algorithme de sélection et de routage}

\subsection{Vue d'ensemble : deux phases}

La résolution du twist 5 s'effectue en deux phases séquentielles et
potentiellement itérables.

\begin{center}
\begin{tikzpicture}[
  box/.style={rectangle, rounded corners=5pt, minimum width=9cm,
              minimum height=1.1cm, align=center, font=\small\bfseries},
  arr/.style={-{Latex[length=3mm]}, thick, gray}
]
  \node[box, fill=bleuléger, draw=theorie] (p0) at (0, 0)
    {Phase 0 — Mise à jour des statuts hospitaliers ($\sigma(h, t)$)};
  \node[box, fill=twistléger, draw=twist] (p1) at (0,-1.7)
    {Phase 1 — Sélection de la destination optimale $h^* \in \mathcal{H}^*(t)$};
  \node[box, fill=vertléger, draw=pratique] (p2) at (0,-3.4)
    {Phase 2 — Calcul de l'itinéraire A*$(s,\, h^*)$ dans le graphe dynamique};
  \node[box, fill=yellow!15, draw=attention] (p3) at (0,-5.1)
    {Phase 3 — Surveillance continue : réseau \emph{et} statuts hospitaliers};
  \node[box, fill=red!8, draw=urgence] (p4) at (0,-6.8)
    {Recalcul si changement détecté (route bloquée OU hôpital saturé)};

  \draw[arr] (p0) -- (p1);
  \draw[arr] (p1) -- (p2);
  \draw[arr] (p2) -- (p3);
  \draw[arr] (p3) -- (p4);
  \draw[arr] (p4.east) -- ++(1.5,0) |- (p0.east)
    node[midway, right, font=\footnotesize, color=urgence]{boucle};
\end{tikzpicture}
\end{center}

\subsection{Pseudocode de la sélection de destination}

\begin{algorithm}[H]
\caption{Sélection du destinataire optimal (Twist 5)}
\begin{algorithmic}[1]
\Require Source $s$, graphe dynamique $G$, ensemble $\mathcal{H}$,
         instant $t$, seuil $\tau$, paramètres $\alpha, \beta$
\Ensure Destination optimale $h^*$ ou \textsc{Escalade}

\State $\mathcal{H}^* \gets \{h \in \mathcal{H} \mid
    \texttt{accepts\_trauma}(h,t) = 1\}$

\If{$\mathcal{H}^* = \emptyset$}
  \State \textbf{déclencher} \textsc{EscaladeZoneComplète}() \Comment{voir §\ref{sec:limites}}
  \State \Return \textsc{null}
\EndIf

\State $\text{meilleur\_coût} \gets +\infty$
\State $h^* \gets \textsc{null}$

\For{$h \in \mathcal{H}^*$}
  \State $d \gets$ A*\_cost$(G, s, h)$ \Comment{temps de trajet estimé}
  \State $\text{pen\_cap} \gets \alpha \cdot \texttt{capacity}(h, t)$
  \State $\text{pen\_seuil} \gets \beta$ \textbf{ si }
         $\texttt{capacity}(h,t) > \tau$ \textbf{ sinon } $0$
  \State $\text{coût} \gets d + \text{pen\_cap} + \text{pen\_seuil}$
  \If{$\text{coût} < \text{meilleur\_coût}$}
    \State $\text{meilleur\_coût} \gets \text{coût}$
    \State $h^* \gets h$
  \EndIf
\EndFor

\State \textbf{journaliser}$(t,\; s,\; h^*,\; \mathcal{H}^*,\;
  \text{meilleur\_coût},\; [\text{coûts de tous les candidats}])$
\State \Return $h^*$
\end{algorithmic}
\end{algorithm}

\subsection{Complexité et performance}

Soit $|\mathcal{H}|$ le nombre total d'hôpitaux dans la zone (typiquement
$3$ à $8$ pour une grande ville africaine). L'algorithme appelle A* une
fois par hôpital candidat. La complexité est donc :

\[
  \mathcal{O}\bigl(|\mathcal{H}^*| \times T_{\text{A*}}\bigr)
  \;\approx\;
  \mathcal{O}\bigl(|\mathcal{H}| \times (|E| + |V| \log |V|)\bigr)
\]

En pratique, avec $|\mathcal{H}| \leq 8$ et un graphe de $\sim$15\,000
nœuds, la sélection complète s'exécute en moins de \textbf{200 ms},
ce qui est compatible avec les contraintes temps réel du système.

\begin{metierbox}{Optimisation : pré-élagage des candidats}
Avant d'appeler A* sur tous les hôpitaux disponibles, le système
élimine d'abord ceux dont la distance à vol d'oiseau dépasse
$2 \times d_{\text{vol}}(s, h^*)$ où $h^*$ est l'hôpital le plus
proche. Cela réduit le nombre d'appels A* à $2$--$3$ en moyenne,
sans risquer d'éliminer un candidat réellement compétitif.
\end{metierbox}

\newpage

% ═══════════════════════════════════════════════════════════════════
\section{Scénario : l'ambulance de Bépanda — suite}

Le scénario d'origine (section 4 du problème) se termine à 14h31 avec
l'ambulance arrivée au CHU. Le twist 5 en reprend le début et modifie
radicalement ce qui se passe entre le départ et la destination.

\subsection{Situation initiale — appel à 14h22}

Il est 14h22. Un enfant de sept ans fait des convulsions à Bépanda.
L'appel arrive au SAMU. Avant de calculer l'itinéraire, le système
interroge les statuts hospitaliers.

\medskip
\begin{center}
\begin{tabular}{lcccc}
  \toprule
  \textbf{Hôpital} & \textbf{Distance estimée} &
  \textbf{Accepte trauma} & \textbf{Taux remplissage} & \textbf{Coût total} \\
  \midrule
  CHU            & 9 min  & \textbf{Non} & 100\% & $+\infty$ \\
  Hôpital Central & 14 min & Oui         & 62\%  & 15,2 min  \\
  Clinique Ngousso & 18 min & Oui         & 41\%  & 18,8 min  \\
  Hôpital de Cité & 22 min & Oui         & 78\%  & 25,6 min  \\
  \bottomrule
\end{tabular}
\end{center}
\medskip

\noindent
Le CHU est immédiatement exclu de $\mathcal{H}^*$ car
$\texttt{accepts\_trauma}(\text{CHU}, 14h22) = 0$.
L'Hôpital Central remporte la sélection avec un coût de 15,2 minutes
(14 min de trajet + $2 \times 0{,}62 = 1{,}24$ min de pénalité capacité,
taux inférieur à $\tau = 0{,}85$ donc pas de pénalité seuil).

\textbf{L'ambulance part à 14h22 en direction de l'Hôpital Central.}

\subsection{Complication en cours de route — 14h27}

Il est 14h27. Deux événements se produisent simultanément :

\begin{enumerate}[leftmargin=2em]
  \item \textbf{Événement routier :} le Boulevard de la Liberté est
    partiellement bloqué (alerte de trafic, délai de propagation :
    4 minutes — situation identique au scénario d'origine).
  \item \textbf{Événement hospitalier :} l'Hôpital Central signale
    une alerte de capacité — son taux de remplissage vient de passer
    à 91\%, et le service trauma est en cours de saturation.
    $\texttt{capacity}(\text{Hôpital Central}, 14h27) = 0{,}91 > \tau$.
\end{enumerate}

Le système recalcule la sélection de destination avec les données
à 14h27 et la position courante de l'ambulance :

\medskip
\begin{center}
\begin{tabular}{lcccc}
  \toprule
  \textbf{Hôpital} & \textbf{Distance restante} &
  \textbf{Accepte trauma} & \textbf{Taux remplissage} & \textbf{Coût total} \\
  \midrule
  CHU            & 8 min  & \textbf{Non} & 100\% & $+\infty$ \\
  Hôpital Central & 9 min  & Oui & 91\% & 15,8 min \\
  Clinique Ngousso & 16 min & Oui & 41\%  & 16,8 min  \\
  Hôpital de Cité & 24 min & Oui & 78\%  & 27,6 min  \\
  \bottomrule
\end{tabular}
\end{center}
\medskip

La pénalité seuil s'applique maintenant à l'Hôpital Central
($0{,}91 > 0{,}85$) : son coût monte à
$9 + 2 \times 0{,}91 + 5 = 15{,}82$ min.
La Clinique Ngousso est maintenant compétitive à 16,8 min.

La différence est faible (1 minute). Le système applique alors un
critère de stabilité : \textbf{il ne change de destination que si le
gain est supérieur à un seuil minimal} $\Delta_{\min} = 2$ minutes,
afin d'éviter des oscillations inutiles.

\begin{metierbox}{Décision à 14h27}
Le gain potentiel en changeant pour la Clinique Ngousso est de
$15{,}82 - 16{,}8 = -0{,}98$ min (Ngousso est plus coûteuse).
L'Hôpital Central reste la meilleure destination.
Le système maintient le cap et note dans le journal la re-validation
de la destination, ainsi que le statut de chaque hôpital au moment
de la vérification.
\end{metierbox}

\subsection{Résolution — 14h36}

L'ambulance arrive à l'Hôpital Central à 14h36. Soit 14 minutes après
le départ — 5 minutes de plus que si le CHU avait été disponible. Le
patient est pris en charge.

\begin{metierbox}{Ce que ce scénario révèle}
Trois décisions ont été prises sur la base d'informations dynamiques :
\begin{enumerate}[leftmargin=1.5em]
  \item à 14h22 — sélection initiale (CHU exclu, Central retenu) ;
  \item à 14h27 — re-validation de la destination malgré la hausse
    de capacité (seuil de stabilité non atteint) ;
  \item à 14h27 — recalcul d'itinéraire suite au blocage routier
    (identique au scénario d'origine, mais vers la nouvelle destination).
\end{enumerate}
Chacune est documentée, reproductible, et explicable.
\end{metierbox}

\newpage

% ═══════════════════════════════════════════════════════════════════
\section{Invariants métier et cas limites}
\label{sec:limites}

\subsection{Invariants introduits par le twist 5}

\begin{invariant}[Non-envoi vers un hôpital saturé]
Le système ne génère jamais d'itinéraire vers un hôpital $h$ tel que
$\texttt{accepts\_trauma}(h, t) = 0$ au moment du calcul.
\end{invariant}

\begin{invariant}[Re-vérification en cours de trajet]
Si un hôpital de destination $h^*$ reçoit une alerte de saturation
alors que le véhicule est en route, le système recalcule la sélection
dans un délai maximal de 30 secondes après réception de l'alerte.
\end{invariant}

\begin{invariant}[Seuil de stabilité]
Un changement de destination en cours de trajet n'est effectué que si
le nouveau coût est inférieur au coût courant d'au moins $\Delta_{\min}$,
afin d'éviter les oscillations contre-productives.
\end{invariant}

\subsection{Le cas critique : aucun hôpital disponible}
\label{subsec:escalade}

Si $\mathcal{H}^*(t) = \emptyset$ — tous les hôpitaux de la zone sont
saturés — le système déclenche une \textbf{procédure d'escalade} :

\begin{enumerate}[leftmargin=2em, itemsep=4pt]
  \item \textbf{Extension géographique} : élargir la recherche à une
    zone de rayon supérieur (par exemple, inclure les hôpitaux
    des villes voisines ou des zones périphériques).
  \item \textbf{Alerte régionale} : notifier le centre de coordination
    régional pour une gestion de crise à plus grande échelle.
  \item \textbf{Destination provisoire} : si la distance vers la
    zone étendue dépasse le délai critique, router vers le poste
    de secours le plus proche (pompiers, caserne) pour stabilisation
    en attendant qu'un hôpital se libère.
  \item \textbf{Journalisation de crise} : l'événement est enregistré
    comme incident critique avec horodatage et état complet du
    réseau hospitalier pour analyse post-incident.
\end{enumerate}

\begin{alertbox}{Cas limite : hôpital se libère pendant le trajet}
Si $\mathcal{H}^*(t) = \emptyset$ au départ mais qu'un hôpital repasse
en mode \texttt{accepts\_trauma = 1} pendant le trajet, le système
annule la procédure d'escalade et recalcule une destination valide.
Ce cas est logué séparément comme « réouverture en cours de trajet ».
\end{alertbox}

\newpage

% ═══════════════════════════════════════════════════════════════════
\section{Implémentation dans le moteur existant}

\subsection{Extension minimale du backend Python}

Le twist 5 nécessite des modifications chirurgicales dans le code
du MVP, sans refonte de l'architecture.

\begin{center}
\begin{tabular}{>{\ttfamily}p{4cm} p{9cm}}
  \toprule
  \textbf{Fichier} & \textbf{Modification} \\
  \midrule
  routing/graph.py &
    Ajouter l'attribut \texttt{hospital\_status} sur les nœuds
    $h \in \mathcal{H}$ ; exposer une méthode
    \texttt{update\_hospital\_status(node\_id, status\_dict)}. \\[0.4em]
  routing/algorithms.py &
    Ajouter la fonction \texttt{select\_destination(G, source, hospitals)}
    qui implémente l'algorithme de sélection (§3). \\[0.4em]
  api/views.py &
    Modifier l'endpoint \texttt{/api/route/} pour accepter une liste
    d'hôpitaux candidats au lieu d'une destination unique ; appeler
    \texttt{select\_destination} avant A*. \\[0.4em]
  api/hospital\_feed.py &
    Nouveau module de polling du statut hospitalier (toutes les 60 s
    via API du ministère de la Santé ou webhook). \\[0.4em]
  routing/logger.py &
    Étendre le schéma de journalisation pour inclure le champ
    \texttt{hospital\_selection\_log}. \\
  \bottomrule
\end{tabular}
\end{center}

\subsection{Schéma de l'appel API modifié}

L'endpoint de routage évolue : au lieu d'une destination fixe,
le client envoie une liste de destinations candidates.

\begin{verbatim}
POST /api/route/
{
  "lat_origin": 3.8921,  "lon_origin": 11.5312,
  "vehicle_type": "ambulance",
  "case_type": "trauma",
  "candidate_hospitals": [
    { "id": "CHU",      "lat": 3.8670, "lon": 11.5178 },
    { "id": "CENTRAL",  "lat": 3.8541, "lon": 11.5023 },
    { "id": "NGOUSSO",  "lat": 3.9012, "lon": 11.5445 }
  ]
}
\end{verbatim}

Le backend filtre les candidats selon leur statut, sélectionne
$h^*$, calcule l'itinéraire et renvoie :

\begin{verbatim}
{
  "selected_hospital": "CENTRAL",
  "selection_reason": "CHU: saturé | CENTRAL: coût 15.2 min | NGOUSSO: 18.8 min",
  "route": { "type": "Feature", "geometry": { ... }, "properties": {
    "cost_minutes": 14.0,
    "total_selection_cost": 15.2
  }},
  "hospital_status_snapshot": {
    "CHU":     { "accepts_trauma": false, "capacity": 1.0 },
    "CENTRAL": { "accepts_trauma": true,  "capacity": 0.62 },
    "NGOUSSO": { "accepts_trauma": true,  "capacity": 0.41 }
  }
}
\end{verbatim}

\newpage

% ═══════════════════════════════════════════════════════════════════
\section{Communication au jury}

\begin{jurybox}{Message clé pour le jury}
\textbf{Le twist 5 déplace le problème de la route vers la destination.}

\medskip
Jusqu'ici, on nous demandait : \og{}Comment arriver \emph{vite} à
l'hôpital ?\fg{} Désormais, la première question est :
\og{}\emph{Quel} hôpital peut encore accueillir ce patient ?\fg{}

\medskip
Notre réponse : la destination devient une variable au même titre
que le trafic. Le système évalue tous les hôpitaux disponibles,
estime le coût réel d'y aller (trajet + risque de saturation à
l'arrivée), choisit le meilleur, et surveille ce choix tout au
long du trajet. Si un hôpital sature en route, il recalcule —
exactement comme il le ferait pour une route bloquée.
\end{jurybox}

\subsection{Questions anticipées du jury}

\paragraph{Question : \og{}Pourquoi ne pas envoyer l'ambulance au CHU
quand même, au cas où ?\fg{}}

\noindent\textbf{Réponse :} Un hôpital qui déclare ne plus accepter
de nouveaux cas de trauma ne peut pas, par définition, prendre en
charge ce patient à l'arrivée. L'envoyer quand même revient à perdre
du temps pour constater un refus sur place — ce qui est catastrophique
pour le patient. La décision de rerouter doit être prise immédiatement,
pas laissée au hasard.

\paragraph{Question : \og{}Comment êtes-vous sûrs que le statut
hospitalier est fiable ?\fg{}}

\noindent\textbf{Réponse :} Comme pour les alertes routières, le
système intègre un score de fiabilité de la source. Un statut
transmis directement par le système informatique de l'hôpital
(intégration API) a une fiabilité haute. Un signalement verbal
par radio a une fiabilité plus faible et déclenche une confirmation
avant de l'utiliser comme critère de rejet. En cas de doute, la
stratégie de repli est de conserver la destination actuelle et
d'envoyer une confirmation humaine.

\paragraph{Question : \og{}Et si tous les hôpitaux sont pleins ?\fg{}}

\noindent\textbf{Réponse :} C'est le cas limite documenté dans
la section~5.2. Le système déclenche une procédure d'escalade :
extension géographique de la recherche, alerte au coordinateur
régional, et si nécessaire, route vers un poste de stabilisation
en attendant qu'une capacité se libère. Ce cas est rare, mais le
système sait comment le gérer — et le journalise comme incident
critique pour analyse ultérieure.

\subsection{Tableau de bord : log de décision hospitalière}

\begin{center}
\begin{tabular}{ll}
  \toprule
  \textbf{Champ du log} & \textbf{Exemple de valeur} \\
  \midrule
  Timestamp de sélection         & 2026-04-26 14:22:07 \\
  Véhicule                       & Ambulance SAMU-07 \\
  Type de cas                    & Trauma pédiatrique \\
  Candidats évalués              & CHU, Hôpital Central, Ngousso, Cité \\
  CHU — statut                   & Refusé (accepts\_trauma = false) \\
  Central — coût sélection       & 15,2 min (retenu) \\
  Ngousso — coût sélection       & 18,8 min \\
  Hôpital retenu                 & Hôpital Central \\
  Itinéraire calculé             & [Bépanda $\to$ Makepe $\to$ Bastos $\to$ Central] \\
  Durée de sélection (algo)      & 87 ms \\
  Re-vérification à 14:27        & Central maintenu (gain Ngousso $< \Delta_{\min}$) \\
  \bottomrule
\end{tabular}
\end{center}

\newpage

% ═══════════════════════════════════════════════════════════════════
\section*{Conclusion}
\addcontentsline{toc}{section}{Conclusion}

Le twist 5 transforme un problème de routage en un problème de
\textbf{planification sous contraintes dynamiques de ressources}.
Ce changement de nature est plus profond qu'il n'y paraît : il oblige
le système à raisonner simultanément sur deux réseaux — le réseau
routier et le réseau de soins — et à optimiser sur les deux à la fois.

\medskip
La solution présentée dans ce document repose sur trois principes
structurants :

\begin{enumerate}[leftmargin=2em, itemsep=4pt]
  \item \textbf{La destination est une variable, pas un paramètre.}
    Le moteur sélectionne la meilleure destination parmi un ensemble
    de candidats avant de calculer l'itinéraire, en intégrant
    disponibilité, capacité et risque de saturation.

  \item \textbf{La surveillance continue couvre désormais le réseau
    de soins.} Le recalcul peut être déclenché par un événement
    hospitalier aussi bien que routier. Les deux types d'événements
    sont traités avec le même mécanisme et la même rigueur.

  \item \textbf{Chaque décision est justifiable.} La sélection de
    destination est entièrement journalisée : hôpitaux évalués, coûts
    calculés, raison du choix, re-vérifications en cours de trajet.
    Si quelqu'un demande compte un mois plus tard, le système peut
    répondre précisément.
\end{enumerate}

\medskip
\begin{center}
  {\color{twist}\rule{0.5\linewidth}{1pt}}\\[0.5em]
  {\itshape \og{}Quand la route change, on recalcule le chemin.\\
  Quand la destination change, on recalcule aussi la destination.\\
  La vraie intelligence, c'est de savoir que les deux peuvent changer en même temps.\fg{}}\\[0.5em]
  {\color{twist}\rule{0.5\linewidth}{1pt}}
\end{center}

\end{document}